\chapter{Directories and the File-System Namespace}\label{ch:directories}
\chapterguide{Chapter 2: files and metadata, especially inode numbers and permissions.}{read absolute and relative paths, distinguish directory entries from inodes, compare hard and symbolic links, and follow recursive traversal in Lab 9.}

\section{Why directories exist}
Storage blocks do not naturally come with human-friendly names such as \texttt{report.txt}. A file system therefore provides a \term{namespace}: a structured way to associate names with file-system objects.

A directory is part of that namespace. In a Unix-like file system, directory entries associate names with underlying objects identified by metadata such as inode numbers.

\section{From flat names to a hierarchy}
Textbooks often describe several directory organizations:
\begin{itemize}
  \item \textbf{single-level:} every file shares one directory;
  \item \textbf{two-level:} each user receives a separate directory;
  \item \textbf{tree-structured:} directories contain files and subdirectories;
  \item \textbf{graph-like structures:} links allow an object to be reachable through more than one path.
\end{itemize}
Modern general-purpose systems use hierarchical trees with linking mechanisms.

\begin{mlfigure}
\begin{tikzpicture}[
  grow=down,
  level distance=1.55cm,
  level 1/.style={sibling distance=4.2cm},
  level 2/.style={sibling distance=2.35cm},
  dir/.style={draw=MLBlue,fill=MLPaleBlue,rounded corners=2mm,minimum width=2.6cm,minimum height=0.75cm,font=\small\bfseries\color{MLNavy}},
  file/.style={draw=MLTeal,fill=MLPaleTeal,rounded corners=1mm,minimum width=1.9cm,minimum height=0.65cm,font=\scriptsize\color{MLNavy}},
  edge from parent/.style={draw=MLMuted,thick,-{Stealth[length=1.8mm]}}
]
\node[dir] {root}
 child {node[dir] {home} child {node[file] {notes.txt}} child {node[file] {code.c}}}
 child {node[dir] {etc} child {node[file] {config}}}
 child {node[dir] {tmp}};
\end{tikzpicture}
\end{mlfigure}

\section{Absolute and relative paths}
An \term{absolute path} begins at the root directory, for example \texttt{/home/student/notes.txt}. A \term{relative path} begins at the process's current working directory.

Two special names appear frequently:
\begin{itemize}
  \item \texttt{.} means the current directory;
  \item \texttt{..} means the parent directory.
\end{itemize}
The shell maintains a current working directory for itself, and a child process normally inherits that starting directory.

\section{Pathname resolution is repeated lookup}
To resolve a path such as \texttt{/home/student/notes.txt}, the system conceptually starts at the root, finds the entry named \texttt{home}, then inside that directory finds \texttt{student}, and finally finds \texttt{notes.txt}.

This is why permissions on intermediate directories matter: reaching a file requires traversing the path, not merely knowing the final filename.

\section{Name, metadata, and data are different layers}
A filename is not the complete file. Metadata can include object type, owner, permissions, timestamps, size, link count, and pointers or mappings to data blocks.

\begin{keyidea}
The namespace tells us \emph{how an object is named}. Metadata tells us \emph{what the object is and its properties}. Storage allocation tells us \emph{where its bytes are represented on storage}.
\end{keyidea}

\section{Hard links}
A hard link gives another directory entry to the same underlying inode. The two names are peers: neither is inherently the ``real'' name. Removing one name does not remove the file data while another hard link still refers to the inode.

The inode's link count helps the file system track how many directory entries refer to it.

\section{Symbolic links}
A symbolic link is a separate small file-system object whose content is a pathname. Following the link causes another pathname lookup.

\begin{mltable}
\renewcommand{\arraystretch}{1.18}
\begin{tabularx}{\linewidth}{@{}
  >{\raggedright\arraybackslash}p{0.28\linewidth}
  >{\raggedright\arraybackslash}X
  >{\raggedright\arraybackslash}X@{}}
\toprule
\rowcolor{UTPPaleBlue!92}
\textbf{Feature} & \textbf{Hard link} & \textbf{Symbolic link} \\
\midrule
Underlying inode & same inode as target & separate inode \\
Stores a pathname & no & yes \\
Can become dangling & not in the same way & yes, if target path disappears \\
Recursive traversal concern & ordinary name to same object & can create path-following cycles \\
\bottomrule
\end{tabularx}
\end{mltable}

\section{Recursive traversal}
A directory walker opens a directory, reads each entry, and recursively repeats the operation for child directories. Two practical rules matter:
\begin{enumerate}
  \item do not recurse into \texttt{.} and \texttt{..};
  \item decide deliberately whether symbolic links will be followed.
\end{enumerate}
Blindly following links can revisit the same locations indefinitely if links form a cycle.

\section{Directory operations and consistency}
Common namespace operations include create, list, rename, link, unlink, and traverse. A real file system must update its metadata consistently even if two processes act at once or a crash interrupts an update.

The lab stays at an observable introductory level: you create a small tree, inspect inode identity and links, and compare the shell's view with a recursive C walker.

\begin{quickcheck}
\begin{enumerate}
  \item What is the starting point of an absolute path? A relative path?
  \item Why can two different names report the same inode number?
  \item Why might a directory walker choose not to follow symbolic links?
\end{enumerate}
\end{quickcheck}

\labsection{9}{File-system namespace, paths, and links}{sec:lab09}
\begin{labproject}{Build a directory tree and investigate names, inode identity, and links}
\textbf{Suggested time:} 90--100 minutes \hfill
\textbf{Level:} introductory/intermediate\par
\textbf{Starter files:} \texttt{make\_demo\_tree.sh}, \texttt{directory\_walk.c}.\par
\textbf{Evidence to keep:} hand-drawn tree, pathname examples, inode/link table, and recursive-walk output.

\begin{labmode}
Use the real Linux/WSL file system. No graphics library is required. The goal is to understand what a directory tree represents, not to draw a decorative visualization.
\end{labmode}

\medskip\noindent\textbf{Learning outcomes.}
\begin{enumerate}
  \item distinguish absolute and relative pathnames;
  \item explain why a filename is not the same thing as the file's metadata/inode;
  \item distinguish a hard link from a symbolic link at an introductory level;
  \item explain how recursive directory traversal follows the namespace.
\end{enumerate}

\medskip\noindent\textbf{Part A --- Create and draw a known hierarchy.}
\begin{lstlisting}[style=osbash]
cd experiments/Lab09
bash make_demo_tree.sh
find demo_tree -maxdepth 4 -print
\end{lstlisting}
Draw the tree by hand. Label directories, regular files, the hard link, and the symbolic link. Your diagram should reflect the names reported by \texttt{find}.

\medskip\noindent\textbf{Part B --- Navigate using different path forms.}
Move into \texttt{demo\_tree/home/student}. Display \texttt{notes.txt} using: (1) a relative path, (2) an absolute path, and (3) a relative path containing \texttt{..}. Run \texttt{realpath notes.txt} and explain what it resolves.

\medskip\noindent\textbf{Part C --- Compare names with inode identity.}
\begin{lstlisting}[style=osbash]
ls -li notes.txt notes-hardlink.txt code-link.c
stat notes.txt
stat notes-hardlink.txt
readlink code-link.c
\end{lstlisting}
Record inode numbers and link counts. Explain why the hard-linked names refer to the same inode while a symbolic link stores a pathname to another object.

\begin{keyidea}
A directory gives names to objects in a namespace. The inode/metadata identifies properties of the underlying file object; storage blocks hold the file's data. These are related layers, not the same thing.
\end{keyidea}

\medskip\noindent\textbf{Part D --- Walk the tree recursively.}
\begin{lstlisting}[style=osbash]
gcc -Wall -Wextra directory_walk.c -o directory_walk
./directory_walk demo_tree
\end{lstlisting}
Find the recursive function call in the source. Explain the stopping rule and why the program deliberately does not recurse through symbolic links.

\medskip\noindent\textbf{Part E --- Controlled namespace change.}
Rename \texttt{notes.txt} while the hard link still exists. Inspect both names again. Then remove one hard link and inspect the remaining link count. Record the observations before deleting the demo tree.

\begin{verification}
Submit the tree, three pathname examples, the inode/link comparison, and a short explanation of why following symbolic links recursively without a policy could create cycles.
\end{verification}
\end{labproject}

